\slajdid{09_3-s032}
\begin{frame}[label=09_3-s032]
\gusttekst\setlength{\abovedisplayskip}{3pt}\setlength{\belowdisplayskip}{3pt}\setlength{\abovedisplayshortskip}{2pt}\setlength{\belowdisplayshortskip}{2pt}PUN mreža kao što smo videli pravi funkciju kod paralelno povezanih tranzistora $\overline A+\overline B$ a kod redno $\overline A\,\overline B$ odnosno najbolje bi bilo da za realizaciju funkcije posmatramo oblik $F$ u kojem se pojavlju promenjive sa komplementnom vrednosti u odnosu na PDN mrežu
\[F=\overline D(A+\overline B\,\overline C)\]
i u tom slučaju svako ILI daje paralelnu vezu a svako I rednu vezu, dok se na ulazima tranzistora nalaze promenljive sa komplementim vrednostima iz ovakvog načina prikazivanja (što odgovara pravim vrednostima iz PDN mreže)
\par\smallskip\centering\input{slike/tikz/s032_pun_slozene_funkcije.tex}
\end{frame}
